		\subsection{主席树（可持久化权值线段树）}

		主席树即\textbf{可持久化权值线段树}，每次单点更新只新建 $O(\log n)$ 个节点，多保存一份「版本根」。处理\textbf{静态区间第 $k$ 小}问题时，第 $i$ 个版本对应前 $i$ 个数构成的权值集合，区间 $[l, r]$ 的第 $k$ 小 $=$ \verb|query(root[l-1], root[r], 1, len, k)|，按左子树差量 \verb|sum[ls[v]] - sum[ls[u]]| 决定向左还是向右走。

		\textbf{由于数组开得很大，本结构体只能放在全局作用域}。

		\begin{minted}{cpp}
struct persisTree {
    // 单点修改区间查询主席树（静态区间第 k 小）
    // 测试 https://www.luogu.com.cn/record/218615977
    // maxn 为点数，maxm 为版本数
    static const int
        maxn = static_cast<int>(5e5) + 7,
        maxm = static_cast<int>(5e5) + 7;
    // 版本数一般不超过 2^20 = 1048576
    static const int LOGmxm = 20;
    static const int
        maxn4 = maxn * 4 + LOGmxm * maxn;
    int ls[maxn4], rs[maxn4], L[maxn4], R[maxn4],
        root[maxm];
    int arr[maxn], sum[maxn4], li[maxn];
    int n = 0, m = 0, idx = 1, Cver = 0, len = 0;

    inline int initLi(int ln, int lm) {
        idx = 1; n = ln; m = lm; Cver = 0;
        root[0] = 1;  // 版本 0 的根是 1
        for (int i = 1; i <= n; ++i) li[i] = arr[i];
        sort(li + 1, li + 1 + n);
        len = unique(li + 1, li + 1 + n) - li - 1;
        return len;
    }
    inline int kth(ll x) {
        return lower_bound(li + 1, li + 1 + len, x)
             - li;
    }
    inline void clone(int node) {
        ++idx;
        L[idx] = L[node];   R[idx] = R[node];
        sum[idx] = sum[node];
        ls[idx] = ls[node]; rs[idx] = rs[node];
    }
    void build(int node, int l, int r) {
        ls[node] = ++idx;  rs[node] = ++idx;
        L[node] = l;       R[node] = r;
        sum[node] = 0;
        if (l == r) return;
        build(ls[node], l, ((l + r) >> 1));
        build(rs[node], (l + r >> 1) + 1, r);
    }
    // node 应填入要 clone 的版本根 root[ver]
    void add(int node, int ver, int p, ll x) {
        clone(node);
        sum[idx] += x;
        if (root[ver] == node) {
            ++Cver;
            root[Cver] = idx;
        }
        if (p == L[node] && p == R[node]) return;
        if (p >= L[ls[node]]
            && p <= R[ls[node]]) {
            int nals = ls[idx];
            ls[idx] = idx + 1;
            add(nals, ver, p, x);
        } else {
            int nars = rs[idx];
            rs[idx] = idx + 1;
            add(nars, ver, p, x);
        }
    }
    // u, v 为版本根；l, r 为值域
    int query(int u, int v, int l, int r, int k) {
        if (l == r) return l;
        int lsum = sum[ls[v]] - sum[ls[u]];
        int mid = l + r >> 1;
        if (lsum >= k)
            return query(ls[u], ls[v], l, mid, k);
        return query(rs[u], rs[v],
                     mid + 1, r, k - lsum);
    }
};
		\end{minted}

		\textbf{典型用法（P3567 KUR-Couriers，求区间是否有出现 $> \lceil \text{len} / 2 \rceil$ 次的数；\href{https://www.luogu.com.cn/record/218724163}{AC 记录}）}：

		\begin{minted}{cpp}
persisTree tr;  // 必须全局

inline void solve() {
    cin >> N >> M;
    tr.init(N, N);
    FOR(i, 1, N) cin >> tr.arr[i];
    int len = tr.initLi();
    tr.build(1, 1, len);
    FOR(i, 1, N) {
        tr.add(tr.root[i - 1], i - 1,
               tr.kth(tr.arr[i]), 1);
    }
    FOR(_, 1, M) {
        ll l, r; cin >> l >> r;
        ll mid = (r - l + 1) / 2 + 1;
        int ans = tr.li[
            tr.query(tr.root[l - 1], tr.root[r],
                     1, len, mid)];
        cout << ans << "\n";
    }
}
		\end{minted}
